%! AP Statistics — Unit 1, Lesson 1.7 — Warm-Up
\documentclass[10pt]{article}
\usepackage{apstats-article}
\usepackage{apstats-boxes}

%=============================================================================
\begin{document}

\pageheader{Unit 1, Lesson 1.7}{Warm-Up}
\namedateperiod

\vspace{0.16in}

\noindent In Lesson 1.6 you described distributions using SOCS --- including
an \textit{informal} estimate of center. Today you will give those estimates
precise numerical values. Use the dotplot below.

\vspace{0.14in}

\begin{center}
\begin{tikzpicture}[scale=1.0]
  % Axis: values 5–50, tick marks every 5 units, scale 0.12 cm/unit
  \draw[thick] (0.5,0) -- (6.2,0);
  \foreach \x/\lbl in {0.60/5, 1.20/10, 1.80/15, 2.40/20, 3.00/25,
                        3.60/30, 4.20/35, 4.80/40, 5.40/45} {
    \draw (\x,-0.12) -- (\x,0.12);
    \node[below, font=\small] at (\x,-0.12) {\lbl};
  }
  \node[below, font=\small\itshape] at (3.0,-0.56)
        {One-way commute time (minutes)};
  % Data: 8, 12, 15, 15, 18, 20, 22, 25, 28, 45
  % x-coord = value * 0.12
  \filldraw[navy] (0.96, 0.38) circle (0.13);   % 8
  \filldraw[navy] (1.44, 0.38) circle (0.13);   % 12
  \filldraw[navy] (1.80, 0.38) circle (0.13);   % 15 (first)
  \filldraw[navy] (1.80, 0.76) circle (0.13);   % 15 (second)
  \filldraw[navy] (2.16, 0.38) circle (0.13);   % 18
  \filldraw[navy] (2.40, 0.38) circle (0.13);   % 20
  \filldraw[navy] (2.64, 0.38) circle (0.13);   % 22
  \filldraw[navy] (3.00, 0.38) circle (0.13);   % 25
  \filldraw[navy] (3.36, 0.38) circle (0.13);   % 28
  \filldraw[navy] (5.40, 0.38) circle (0.13);   % 45
\end{tikzpicture}
\end{center}

\vspace{0.18in}

\begin{enumerate}[leftmargin=1.5em, itemsep=10pt]

  \item How many students are represented in this dotplot? What is the
        variable being measured?\\[6pt]
        \writeline

  \item What is the \textbf{range} of commute times? Show your calculation.\\[6pt]
        \writeline

  \item List the data values in order from least to greatest. Then identify
        the \textbf{middle value} (or middle two values) of the sorted list.\\[6pt]
        \writeline\\[1pt]\writeline

  \item One commute time appears to be much larger than the rest. Which value
        is it, and how does it stand out in the dotplot?\\[6pt]
        \writeline\\[1pt]\writeline

\end{enumerate}

\vspace{0.14in}

\begin{tcolorbox}[colback=greenbg, colframe=greenacc, boxrule=0.8pt, arc=2mm,
    left=4mm, right=4mm, top=3mm, bottom=3mm]
\small
\begin{itemize}[leftmargin=1.3em, itemsep=5pt]
  \item The \textbf{median} is exactly that middle value you identified above ---
        the precise measure of center that resists extreme values.
  \item Today we also calculate the \textbf{mean}, \textbf{IQR}, and
        \textbf{standard deviation}, and learn when each is the right choice.
  \item The large value at 45 minutes will motivate the \textbf{outlier rule}.
\end{itemize}
\end{tcolorbox}

\vspace{0.12in}

\noindent\textcolor{slate}{\small One question you have going into today's lesson,
or something you're not sure about yet:}\\[8pt]
\writeline

\end{document}
